\section{Conclusion}\label{sec:conclusion}

We set out to measure what changes when an AI agent asks about wearable data through a broker that computes the
answer, releases only what the question needs, and records every decision. Holding the agent model fixed, \sys{}
raised accuracy for the models that were weak with raw context (\nSigCA{} of \nTestedCA{} at $p<0.05$ against raw
context; pooled \bootPooledCA{} points) and lowered it significantly for none, though it was worse on a sleep
threshold question where its planner failed. For the strongest models no accuracy difference was detectable; for
them the benefit was a shorter context, a release limited to the granted level and an audit record, at the cost of
an extra planner call. Raw context over-exposed the record by a factor of \expRatioAC{}, and under a corrected,
hand-audited leak detector \sys{} leaked in \leakPooledC{} of privacy runs against \leakPooledA{} for raw context and
\leakPooledB{} for open tools. The remaining errors sit in planning and in the agent model's reading of the packet,
where the broker's logged reasons make them visible; one round of fixes to the planner is measured in the second
iteration. Code, data and the scripts that produce every number in this paper are available at \codeurl{}.
